% !TEX root = ../../main.tex

\section{本章小结}

本章围绕 HearBridge 系统的实现需求，对相关技术进行了分析。
HarmonyOS、ArkTS 和 ArkUI 为移动端手语资源展示、训练交互和结果展示提供了基础；
Spring Boot、MyBatis、MySQL、Redis 和 MinIO 支撑了后端业务接口、数据持久化、登录态管理和文件资源存储；
Vue 3、TypeScript、Vite 和 Element Plus 支撑了后台管理端的页面开发和资源维护；
FastAPI、MediaPipe 和 TensorFlow/Keras 构成了 Python 手势识别服务的核心技术基础。

此外，本章还分析了 HTTP、WebSocket 和 DeepSeek 等关键支撑技术。
总体来看，本文采用的技术路线能够较好地满足系统在移动端训练、后台资源管理、识别服务调用和句子视频识别方面的需求。
下一章将进一步从用户角色、业务流程和非功能性要求等角度对 HearBridge 系统进行需求分析。
